\documentclass[twoside,12pt]{book}
\usepackage{graphicx}
\usepackage{amsmath}
\graphicspath{{/images}}
%\usepackage[export]{adjustbox}
\usepackage[table,xcdraw]{xcolor} % Más atrás da incompatibilidades con otros paquetes


%%%%%%%%%%%%%%%%%%%%%%%%%%%%%%%%%%%%%%%%%%%%%%%%%%%%%%%%%%%%%%
%% Estilo TFG UC3M
\usepackage{tfg_uc3m}

%%%%%%%%%%%%%%%%%%%%%%%%%%%%%%%%%%%%%%%%%%%%%%%%%%%%%%%%%%%%%%
%% Otros paquetes que quieras usar
\usepackage{listings}
\usepackage{svg} % Para insertar gráficos vectoriales SVG
\usepackage{lipsum} % Para mostrar "Lorem ipsum"

% Para poder utilizar otras fuentes (requiere un compilador como LuaLaTeX)
\usepackage{fontspec}
\setmainfont{Arial}
%\setmainfont{Times New Roman}

% Para trabajar fácilmente con emojis (requiere un compilador como LuaLaTeX)
\usepackage{emoji}

%%%%%%%%%%%%%%%%%%%%%%%%%%%%%%%%%%%%%%%%%%%%%%%%%%%%%%%%%%%%%%
%\usepackage{biblatex}

\custombib{\bibtype}
	
\addbibresource{bibliography/bibliography.bib} %% ficheros con la bibliografía


\definecolor{mygreen}{rgb}{0,0.6,0}
\definecolor{mygray}{rgb}{0.47,0.47,0.33}
\definecolor{myorange}{rgb}{0.8,0.4,0}
\definecolor{mywhite}{rgb}{0.98,0.98,0.98}
\definecolor{myblue}{rgb}{0.01,0.61,0.98}

\lstset{ %
  backgroundcolor=\color{mywhite},   
  basicstyle=\footnotesize,       
  breakatwhitespace=false,         
  breaklines=true,                 
  captionpos=b,                   
  commentstyle=\color{mygray},    
  deletekeywords={...},           
  escapeinside={\%*}{*)},          
  extendedchars=true,              
  frame=shadowbox,                    
  keepspaces=true,                 
  keywordstyle=\color{myorange},       
  language=Octave,                
  morekeywords={*,...},            
  numbers=left,                    
  numbersep=5pt,                   
  numberstyle=\tiny\color{mygray}, 
  rulecolor=\color{black},         
  rulesepcolor=\color{myblue},
  showspaces=false,                
  showstringspaces=false,          
  showtabs=false,                  
  stepnumber=2,                    
  stringstyle=\color{myorange},    
  tabsize=2,                       
  title=\lstname                   
}    

\usepackage{tikz}
\def\checkmark{\tikz\fill[scale=0.4](0,.35)--(.25,0)--(1,.7)--(.25,.15)--cycle;}

%%%%%%%%%%%%%%%%%%%%%%%%%%%%%%%%%%%%%%%%%%%%%%%%%%%%%%%%%%%%%%
%% Glosario
\usepackage[acronym,toc]{glossaries}
\makeglossaries

\loadglsentries{others/glossary}
%%%%%%%%%%%%%%%%%%%%%%%%%%%%%%%%%%%%%%%%%%%%%%%%%%%%%%%%%%%%%%
%% Empieza el documento
\begin{document}
\selectlanguage{english}
\pagestyle{empty}

%%%%%%%%%%%%%%%%%%%%%%%%%%%%%%%%%%%%%%%%%%%%%%%%%%%%%%%%%%%%%%
%% Portada
\makecover

\cleardoublepage

%%%%%%%%%%%%%%%%%%%%%%%%%%%%%%%%%%%%%%%%%%%%%%%%%%%%%%%%%%%%%%
%% Dedicatoria

\hfill{\itshape \textbf{Acknowledgments.}}
\vspace{1.2cm}

\hfill{\itshape This is an example.}
\vspace{1cm}

\hfill{\itshape And this is another example.}
\vspace{1cm}
\clearpage

\pagenumbering{Roman} % La numeración debe ser romana hasta el primer capítulo
\pagestyle{plain}  % Debe ir junto a la numeración para empezar a contar las páginas

%%%%%%%%%%%%%%%%%%%%%%%%%%%%%%%%%%%%%%%%%%%%%%%%%%%%%%%%%%%%%%
%% Abstract in English
\begin{otherlanguage}{english}

\chapter*{Abstract}

% Text body
\lipsum[1-2]

\vspace{1cm}

\textbf{Keywords:}

% Keywords
Term 1; term 2; ...; Term n
	 
\vfill


\end{otherlanguage}
\clearpage

%%%%%%%%%%%%%%%%%%%%%%%%%%%%%%%%%%%%%%%%%%%%%%%%%%%%%%%%%%%%%%
%% Índice de contenidos
\tableofcontents
\clearpage

%%%%%%%%%%%%%%%%%%%%%%%%%%%%%%%%%%%%%%%%%%%%%%%%%%%%%%%%%%%%%%
%% Otros índices
\listoffigures
\clearpage

\listoftables
\clearpage

%\listequationsname
%\cleardoublepage

%%%%%%%%%%%%%%%%%%%%%%%%%%%%%%%%%%%%%%%%%%%%%%%%%%%%%%%%%%%%%%
%% 
\pagenumbering{arabic} % Iniciamos la numeración árabe en el primer capítulo
\pagestyle{fancy}
% Remembering the chapter title:
\renewcommand{\chaptermark}[1]{\markboth{\headertitlefont\thechapter\quad#1}{}}
% Remembering section number and title:
\renewcommand{\sectionmark}[1]{\markright{\headertitlefont\thesection\quad#1}}
% Clear all header fields
\fancyhead{}
\fancyhead[RE]{}{}
\fancyhead[LE]{\leftmark}
\fancyhead[RO]{\rightmark}
\setlength{\headheight}{14.5pt}
% Clear all footer fields
\fancyfoot{}
\fancyfoot[C]{\thepage}
\renewcommand{\headrulewidth}{0pt}
\renewcommand{\footrulewidth}{0pt}

%%%%%%%%%%%%%%%%%%%%%%%%%%%%%%%%%%%%%%%%%%%%%%%%%%%%%%%%%%%%%%
%% Capítulos

%% Capítulo inicial a comentar cuando se escriba el documento final
%% Ejemplo de comandos en LaTeX
\chapter{Elementos básicos}
\label{chap:básicos} % Etiqueta para enlazar al capítulo desde el texto

% Definición de colores personalizados
\definecolor{azulUC3M}{RGB}{24, 30, 146}
\definecolor{amarilloFosforito}{HTML}{FFFF00}


\section{Texto}

\subsection{Texto simple}

Mi primer documento en \LaTeX. Ejemplo simple
sin paquetes adicionales.

% -------------------------

\subsubsection{Formatos para texto simple}

Lo más \textbf{importante}
de este \underline{documento} 
no es la \textit{cursiva}.
También podemos \textbf{\textit{encadenar formatos}}.

% -------------------------

\subsection{Colores en el texto}

Hola, soy un \textcolor{red}{texto en rojo} y yo un \textcolor{blue}{texto en azul} y en \textcolor{green}{verde}.

% -------------------------

\section{Texto con color personalizado y subrayado de color}

Ahora podemos poner \textcolor{azulUC3M}{UC3M} en color corporativo y \colorbox{amarilloFosforito}{subrayar texto} en amarillo ``fosforito’’.

% -------------------------

\section{Salto de página}

% Salto de página
\newpage

% -------------------------

\section{Salto de línea}

Texto
% Salto de línea
\newline
Otro texto en nueva línea

% -------------------------

\section{Espaciado personalizado entre líneas de texto}

Hola, soy un texto

% Espaciado personalizado
\vspace{1cm}
Y yo soy otro texto con una separación de 1 cm.

\vspace{50pt}
Y yo soy otro texto más alejado.




\section{Listas simples}

Las listas son sencillas. Para crearlas:
\begin{itemize}
  \item Las entradas de la lista comienzan con el comando \verb|\item| .
  \item Este tipo de lista tiene cada entrada empezando con un ``bullet''.
  \item La longitud de cada entrada es arbitraria.
  \item Y el número de elementos también lo es.
\end{itemize}

% -------------------------

\section{Listas enumeradas}

Las listas enumeradas también son sencillas:
\begin{enumerate}
  \item Los elementos se enumeran automáticamente.
  \item Los números  se crean automáticamente en el entorno \texttt{enumerate}.
  \item Otra entrada
\end{enumerate}

% -------------------------

\subsection{Sublistas enumeradas}

Ejemplo de sublistas enumeradas:
\begin{enumerate}
  \item Primer nivel.
  \begin{enumerate}
    \item Segundo nivel.
    \begin{enumerate}
      \item Tercer nivel.
      \item Tercer nivel.
    \end{enumerate}
    \item Segundo nivel.
  \end{enumerate}
  \item Primer nivel.
\end{enumerate}

% -------------------------

\section{Más sublistas}

Otro ejemplo de sublistas:
\begin{itemize}
  \item Primer nivel.
  \begin{itemize}
    \item Segundo nivel.
    \begin{itemize}
      \item Tercer nivel.
      \item Tercer nivel.
    \end{itemize}
    \item Segundo nivel.
  \end{itemize}
  \item Primer nivel.
\end{itemize}

% -------------------------

\section{Listas personalizadas}

Tanto las listas con puntos como las enumeradas permiten cambiar el elemento de inicio con \verb|\item[elemento]|
\begin{enumerate}
  \item Primer punto
  \item Segundo punto
  \item[!] Tercer punto, comenzando con exclamación
  \item[$\blacksquare$] Ahora empezando con un cuadrado negro
  \item[NOTA] Esta entrada no tiene ``bullet‘’
  \item[] Entrada en blanco?
\end{enumerate}



% -------------------------
% -------------------------
% -------------------------



\section{Imágenes}

\subsection{Raster}

\begin{figure}[ht]
    \centering
    \includegraphics[height=7cm]{images/examples/gul_logo_raster.png}
    \caption{Logo del GUL en formato PNG}
    \label{fig:logo1}
\end{figure}

La figura \ref{fig:logo1} se encuentra en la página \pageref{fig:logo1}.

\newpage

% -------------------------

\subsection{Vectorial (svg)}

\begin{figure}[ht]
     \includesvg[width=\textwidth]{images/examples/gul_logo_vector.svg}
     \centering
     \caption{Logo del GUL en formato vectorial (svg)}
     \label{fig:logo2}
\end{figure}



\section{Tablas y enlaces externos (links)}

% ------------------------------------------------------
% Recurso recomendado: https://www.tablesgenerator.com/
% ------------------------------------------------------
Recurso recomendado: \href{https://www.tablesgenerator.com/}{Tables generator}.

\begin{table}[ht]
  \begin{tabular}{ |c|c|c| } 
    \hline
    Columna A & Columna B & Columna C\\ 
    1 & 3 & 5 \\ 
    2 & 4 & 6 \\ 
    \hline
  \end{tabular}
  \centering
  \caption{Tabla de ejemplo}
  \label{tab:tabla1}
\end{table}

\vspace{1cm}

\begin{table}[ht]
\begin{tabular}{ccccc}
\cline{1-4}
\multicolumn{2}{|c|}{emosido}                       & \multicolumn{1}{c|}{}      & \multicolumn{1}{c|}{}     &                       \\ \cline{1-4}
\multicolumn{1}{|c|}{} & \multicolumn{2}{c|}{\multirow{2}{*}{\textbf{engañado}}} & \multicolumn{1}{c|}{}     &                       \\ \cline{1-1} \cline{4-4}
\multicolumn{1}{|c|}{} & \multicolumn{2}{c|}{}                                   & \multicolumn{1}{c|}{}     &                       \\ \cline{1-4}
\multicolumn{1}{|c|}{} & \multicolumn{1}{c|}{}      & \multicolumn{1}{c|}{}      & \multicolumn{1}{c|}{hola} &                       \\ \cline{1-4}
                       &                            &                            &                           & \textit{Hello there!}
\end{tabular}
\centering
\caption{Otra tabla de ejemplo}
\label{tab:tablaEmosido}
\end{table}



\section{Emojis}

% -------------------------

\subsection{Emojis con el paquete ``emoji''}

Emojis con colores gracias al paquete \texttt{emoji} y LuaLaTeX:
\emoji{leaves}
\emoji{rose}

Otro ejemplo, en blanco y negro, usando el código Unicode para carita feliz (U+1F600): {\fontspec{Symbola}\symbol{"1F600}}

% -------------------------

\subsection{Emojis con el tipografía específica (NotoColorEmoji.ttf) y códigos de caracter Unicode}

\newfontfamily\emojifont[Renderer=Harfbuzz]{NotoColorEmoji.ttf}

Emojis utilizando una tipografía específica y código de caracter Unicode (U+1F600):
% Utilizar \emojifont en un grupo para hacer que su efecto sea local
{\emojifont \Uchar"1F600 \char"1F600}
% http://www.unicode.org/emoji/charts/full-emoji-list.html



\section{Acrónimos y glosario}

El \gls{GULa} (o \gls{GULg}) forma parte de la \gls{UC3Ma} (\gls{UC3Mg}).



\section{Bibliografía y citas}

Cita a una referencia \cite{prieto2014applications}, y cita con dos referencias \cite{prieto2014applications, gonzalez2020simultaneous}.

\newpage


\chapter{Introduction}
\label{chap:01intro}

Write here a small intro to this chapter.

\section{Context}

\section{Problem}

\section{Motivation}

\section{Objective}

\section{Regulatory framework}

\section{Structure of the document}

\newpage

\chapter{State of the art}
\label{chap:02soa}

Write here a small intro to this chapter.

\section{Current situation}

\section{Existing framework/solutions}

\section{Comparative chart among solutions}

\newpage

\chapter{Proposed solution / methods}
\label{chap:03sol-met}

\newpage

\chapter{Validation and testing}
\label{chap:04valtest}

\newpage

\chapter{Results and discussion}
\label{chap:05res-disc}


\newpage

\chapter{Project management}
\label{chap:06mgmnt}

\newpage

\chapter{Conclusions}
\label{chap:07conc}

\newpage

\chapter{Future research lines}
\label{chap:08future}

\newpage

\printglossary[type=\acronymtype,title=Acronyms]
\printglossary[type=main]

%%%%%%%%%%%%%%%%%%%%%%%%%%%%%%%%%%%%%%%%%%%%%%%%%%%%%%%%%%%%%%
%% Bibliografía
\printbibliography

%%%%%%%%%%%%%%%%%%%%%%%%%%%%%%%%%%%%%%%%%%%%%%%%%%%%%%%%%%%%%%
%% Anexos

% \input{capitulos/08anexos}

\cleardoublepage


\end{document}
